\documentclass[11pt]{article}
\usepackage[localtheorems]{raeez-math-template}

\newtheorem{theorem}{Theorem}
\newtheorem{proposition}[theorem]{Proposition}
\newtheorem{lemma}[theorem]{Lemma}
\newtheorem{corollary}[theorem]{Corollary}
\theoremstyle{definition}
\newtheorem{definition}[theorem]{Definition}
\newtheorem{construction}[theorem]{Construction}
\theoremstyle{remark}
\newtheorem{remark}[theorem]{Remark}

\newcommand{\CC}{\mathbb{C}}
\newcommand{\ZZ}{\mathbb{Z}}
\newcommand{\HH}{\mathbb{H}}
\newcommand{\cA}{\mathcal{A}}
\newcommand{\cH}{\mathcal{H}}
\newcommand{\cO}{\mathcal{O}}
\newcommand{\Db}{D^b}
\newcommand{\Coh}{\mathrm{Coh}}
\newcommand{\Tr}{\mathrm{Tr}}
\newcommand{\str}{\mathrm{str}}
\newcommand{\SL}{\mathrm{SL}}
\newcommand{\ch}{\mathrm{ch}}
\newcommand{\Irr}{\mathrm{Irr}}
\newcommand{\phiEZ}{\phi_{0,1}^{\mathrm{EZ}}}

\title{The K3 Elliptic Genus, Mathieu Moonshine, and the \texorpdfstring{$\Delta_5$}{Delta5} Normalisation}
\author{Raeez Lorgat}
\date{2026-04-22}

\begin{document}
\maketitle

\begin{abstract}
The K3 elliptic genus is the weak Jacobi form
\[
  Z_{K3}(\tau,z)=2\,\phiEZ(\tau,z),
  \qquad
  \phiEZ\big|_{q^0}=y^{-1}+10+y.
\]
Its small \(\mathcal N=4\) character decomposition isolates the
Eguchi-Ooguri-Tachikawa mock modular form
\[
  H(\tau)=2q^{-1/8}\bigl(-1+45q+231q^2+770q^3+\cdots\bigr),
\]
whose Zwegers completion has shadow line spanned by \(\eta^3\).  In
the EOT/DMZ Eichler-integral normalisation this completion is written
with the \(24\,\eta^3\) shadow convention.
The same \(\phiEZ\) is the primitive Jacobi input for the
Gritsenko-Nikulin denominator \(D_5=64^{-1}\Delta_5\), giving
\(\kappa_{\mathrm{BKM}}(\Delta_5)=c_1(0)/2=5\).  The Igusa form
\(\Phi_{10}\) belongs to the square normalisation lane
\(\Phi_{10}^{\mathrm{un}}=\Delta_5^2\).
\end{abstract}

\tableofcontents

\section{Jacobi Normalisation}
\label{sec:jacobi-normalisation}

Let \(q=e^{2\pi i\tau}\) and \(y=e^{2\pi iz}\).  The Eichler-Zagier
normalisation of the weight \(0\), index \(1\) weak Jacobi form is
\[
  \phiEZ(\tau,z)=y^{-1}+10+y+O(q).
\]
Equivalently, if
\[
  \phiEZ(\tau,z)=\sum_{n\geq 0}\sum_{\ell\in\ZZ}
       c(4n-\ell^2)\,q^n y^\ell,
\]
then
\[
  c(-1)=1,\qquad c(0)=10,\qquad
  \phiEZ(\tau,0)=12.
\]
The K3 elliptic genus is
\begin{equation}
\label{eq:k3-elliptic-genus}
  Z_{K3}(\tau,z)=2\,\phiEZ(\tau,z).
\end{equation}
Its head is therefore
\[
  Z_{K3}\big|_{q^0}=2y^{-1}+20+2y,
  \qquad Z_{K3}(\tau,0)=24=\chi_{\mathrm{top}}(K3).
\]

\begin{remark}[Four invariants]
\label{rem:four-invariants}
The coefficient \(2\) in \eqref{eq:k3-elliptic-genus} is the K3
chiral and categorical value
\[
  \kappa_{\mathrm{ch}}(V_{K3})
  =\kappa_{\mathrm{cat}}(K3)
  =\chi(\cO_{K3})=2.
\]
It is distinct from the topological Euler number \(24\).  On
\(K3\times E\) the four structural values remain separated:
\[
  \kappa_{\mathrm{cat}}(K3\times E)=0,\qquad
  \kappa_{\mathrm{ch}}^{\mathrm{Heis}}(K3\times E)=3,\qquad
  \kappa_{\mathrm{BKM}}(\Delta_5)=5,\qquad
  \kappa_{\mathrm{fiber}}(K3\times E)=24.
\]
Each value is attached to its construction.
\end{remark}

\section{The \texorpdfstring{$\mathcal N=4$}{N=4} Decomposition}
\label{sec:n4-decomposition}

The Ramond-sector elliptic genus of a K3 sigma model is
\[
  Z_{K3}(\tau,z)
  =\str_{\cH_{\mathrm{RR}}}
     (-1)^F y^{J_0^3}q^{L_0-c/24}.
\]
At central charge \(c=6\), the small \(\mathcal N=4\) algebra has
short characters and long characters.  Eguchi-Ooguri-Tachikawa write
\begin{equation}
\label{eq:eot-decomposition}
  Z_{K3}(\tau,z)
  =24\,\mu(\tau,z)\,
       \frac{\vartheta_1(\tau,z)^2}{\eta(\tau)^3}
   +H(\tau)\,
       \frac{\vartheta_1(\tau,z)^2}{\eta(\tau)^3},
\end{equation}
where
\[
  \mu(\tau,z)=
  \frac{-i\,y^{1/2}}{\vartheta_1(\tau,z)}
  \sum_{n\in\ZZ}
    \frac{(-1)^n q^{n(n+1)/2}y^n}{1-q^ny}
\]
is the Zwegers Appell-Lerch sum and
\begin{equation}
\label{eq:eot-H}
  H(\tau)=2q^{-1/8}
   \bigl(-1+\sum_{n\geq 1}A_nq^n\bigr).
\end{equation}
The first coefficients are
\[
\begin{array}{c|rrrrrrrrrr}
n&1&2&3&4&5&6&7&8&9&10\\
\hline
A_n&45&231&770&2277&5796&13915&30843&65550&132825&258760.
\end{array}
\]

\begin{proposition}[Polar and head constants]
\label{prop:polar-head}
The polar term of \(H\) is \(-2q^{-1/8}\), and the massless
coefficient in \eqref{eq:eot-decomposition} is \(24\).
\end{proposition}

\begin{proof}
Equation~\eqref{eq:eot-H} gives the polar coefficient
\(2\cdot(-1)=-2\), equal to
\(-\kappa_{\mathrm{ch}}(V_{K3})\).  Equation~\eqref{eq:eot-decomposition}
places \(24=\chi_{\mathrm{top}}(K3)\) in front of the Appell-Lerch
short-character contribution.  The two constants come from different
parts of the \(\mathcal N=4\) character expansion.
\end{proof}

\section{Completion and Shadow}
\label{sec:completion-shadow}

The Appell-Lerch sum \(\mu\) has a non-holomorphic Zwegers completion.
Consequently \(H\) is the holomorphic part of a harmonic Maass form of
weight \(1/2\).  Fix the EOT/DMZ Eichler-integral convention by
\[
  \widehat H_{\mathrm{EOT}}(\tau)
  =
  H(\tau)+24\,\mathcal E_{\eta^3}(\tau),
\]
where \(\mathcal E_{\eta^3}\) is the Eichler integral of
\(\eta(\tau)^3\) in the DMZ kernel.  This convention is recorded by
\[
  S_{\mathrm{EOT}}(\tau)=24\,\eta(\tau)^3.
\]
With the Bruinier-Funke operator
\(\xi_{1/2}=2iy^{1/2}\overline{\partial_{\bar\tau}}\), applying
\(\xi_{1/2}\) to \(\widehat H_{\mathrm{EOT}}\) recovers a fixed non-zero
scalar multiple of \(24\,\eta^3\).  The scalar is determined by the
normalisation of the Eichler kernel and of \(\xi_{1/2}\).  The invariant
modular object is the one-dimensional shadow line
\[
  \CC\cdot\eta^3 \subset M_{3/2}.
\]
The coefficient \(24\) belongs to the short-character trace in
\eqref{eq:eot-decomposition} and to the EOT/DMZ completion convention.
The Borcherds constants \(c_N(0)\) belong to the denominator-product
construction of Section~\ref{sec:delta5}.

\begin{remark}[Two shadow languages]
\label{rem:two-shadow-languages}
The analytic shadow line is \(\CC\cdot\eta^3\) in weight \(3/2\).  The
shadow obstruction tower of a chiral algebra is a bar-complex tower of
obstructions.  For \(V_{K3}\) the two languages meet through the
\(\mathcal N=4\) trace and its Appell-Lerch completion.  The asserted
equality is the analytic EOT/DMZ shadow statement.
\end{remark}

\section{Mathieu Characters}
\label{sec:mathieu-characters}

The EOT observation begins with half-coefficients:
\[
  45,\quad 231,\quad 770,\quad 2277,\quad 5796,\ldots
\]
which become \(M_{24}\)-representation dimensions after the prefactor
\(2\) in \eqref{eq:eot-H} is restored.  The character-theoretic
statement uses twining.  For \(g\in M_{24}\),
Gaberdiel-Hohenegger-Volpato define twined genera by replacing
\(\dim H_n\) with \(\Tr_{H_n}(g)\) in the \(\mathcal N=4\) expansion:
\[
  \phi_g(\tau,z)
  =
  \frac12\left(
    \Tr_{H_{00}}(g)\,\ch_{1/4,0}^{\mathcal N=4}(\tau,z)
    -
    \Tr_{H_0}(g)\,\ch_{1/4,1/2}^{\mathcal N=4}(\tau,z)
    +
    \sum_{n\geq 1}\Tr_{H_n}(g)\,
       \ch_{n+1/4,1/2}^{\mathcal N=4}(\tau,z)
  \right).
\]
At the identity class,
\[
  A_n=\frac12\Tr_{H_n}(1),\qquad n\geq 1.
\]
For an element of order \(N\), \(\phi_g\) is a weight \(0\), index \(1\)
Jacobi form for \(\Gamma_0(N)\), with the multiplier determined by the
GHV table.  The ordinary \(M_{24}\) character table has \(26\) conjugacy
classes and \(26\) irreducible characters.  Compressed twining tables
and ordinary character tables are separate objects: power-conjugate
classes can have the same twined Jacobi form, while character
reconstruction uses all ordinary character-table columns.

\begin{theorem}[Gannon]
\label{thm:gannon}
For every \(n\geq 1\), the class function
\[
  g\longmapsto \Tr_{H_n}(g)
\]
is the character of a genuine finite-dimensional \(M_{24}\)-representation.
The finite boundary terms \(H_{00}\) and \(H_0\) are fixed by the
massless and polar \(\mathcal N=4\) characters; the minus sign before
\(H_0\) is part of the character decomposition.
\end{theorem}

\begin{remark}[Dimensions and decompositions]
\label{rem:dimension-vs-character}
The single integer \(A_n=\frac12\Tr_{H_n}(1)\) records only the
identity-class half-trace in the EOT normalisation.  Many integer sums
of irreducible dimensions can share that value.  The full class
function \(g\mapsto \Tr_{H_n}(g)\) determines the irreducible
multiplicities by the usual orthogonality of characters.
\end{remark}

\section{\texorpdfstring{$\Delta_5$}{Delta5} and the Square Lane}
\label{sec:delta5}

The same Jacobi form \(\phiEZ\) controls the Borcherds denominator
attached to \(K3\times E\).  In the Eichler-Zagier convention the
constant coefficient is
\[
  c_1(0)=10.
\]
The Borcherds weight formula gives
\begin{equation}
\label{eq:bkm-weight-delta5-polyakov}
  \kappa_{\mathrm{BKM}}(\Delta_5)
  =
  \mathrm{wt}(\Delta_5)
  =
  \frac{c_1(0)}{2}
  =
  5.
\end{equation}
For the primitive CHL denominator ladder the constants are
\[
\begin{array}{c|ccccc}
N&1&2&3&4&6\\
\hline
c_N(0)&10&8&6&4&2\\
\kappa_{\mathrm{BKM}}(\Phi_N)=c_N(0)/2&5&4&3&2&1
\end{array}
\]
This is the CHL/Gritsenko-Nikulin denominator scope.  Ordinary
Mathieu twining genera supply \(M_{24}\) class functions; a class of
order \(N\) enters the displayed constants only after the primitive CHL
quotient datum and its Borcherds input are fixed.

The theta-normalised form satisfies
\[
  [q^{1/2}r^{1/2}s^{1/2}]\Delta_5=64.
\]
The monic denominator used in the Oberdieck-Pandharipande scalar
normalisation is
\begin{equation}
\label{eq:D5-normalisation}
  D_5=64^{-1}\Delta_5.
\end{equation}
The denominator algebra uses \(D_5(2Z)=64^{-1}\Delta_5(2Z)\).
The square normalisations are
\[
  \Phi_{10}^{\mathrm{un}}=\Delta_5^2,
  \qquad
  \Phi_{10}^{\mathrm{OP}}=D_5^2=4096^{-1}\Delta_5^2.
\]
Thus \(\Delta_5\) is the primitive denominator in the
\(\phiEZ\)-Borcherds lane, while \(\Phi_{10}\) records the square
normalisation data.

\section{CY-A\texorpdfstring{$_2$}{2} Reading}
\label{sec:cy-a2-reading}

For a K3 surface \(S\), the category \(\Db\Coh(S)\) is Calabi-Yau of
dimension \(2\).  The CY-to-chiral construction at \(d=2\) has
\(E_2\)-chiral scope.  At the object level, after choosing the
sigma-model \(\mathcal N=4\) specialisation, the chiral output has genus
one trace
\[
  \Tr_{V_{K3}}
  \bigl((-1)^F y^{J_0^3}q^{L_0-c/24}\bigr)
  =
  Z_{K3}(\tau,z).
\]
The mock modular form \(H\) is therefore the massive \(\mathcal N=4\)
summand of this \(d=2\) trace.  The statement uses the \(d=2\)
\(E_2\)-chiral output and its stage-two sigma-model specialisation.
All proof data for \eqref{eq:eot-decomposition} live in this \(d=2\)
trace calculation.

\begin{theorem}[Standalone K3 mock modular statement]
\label{thm:standalone-k3-mock}
Let \(V_{K3}\) be the small \(\mathcal N=4\) K3 sigma-model vertex
algebra at \(c=6\).  The elliptic genus of \(V_{K3}\) is
\(Z_{K3}=2\phiEZ\).  Its \(\mathcal N=4\) character decomposition
contains the holomorphic function \(H\) of \eqref{eq:eot-H}.  The
completion \(\widehat H\) is a harmonic Maass form of weight \(1/2\)
whose EOT/DMZ completion convention is governed by \(24\,\eta^3\) and
whose operator-normalised shadow line is \(\CC\cdot\eta^3\).  For
\(n\geq 1\), the coefficient class functions of the corresponding
twined forms are genuine \(M_{24}\)-characters in the sense of Gannon.
\end{theorem}

\begin{proof}
The identity \(Z_{K3}=2\phiEZ\) follows from the Eichler-Zagier
classification of \(J_{0,1}^{\mathrm{weak}}\) and the K3 head
\(2y^{-1}+20+2y\).  Eguchi-Ooguri-Tachikawa decompose this Jacobi
form as \eqref{eq:eot-decomposition}, with \(H\) given by
\eqref{eq:eot-H}.  Zwegers completion of the Appell-Lerch sum supplies
the Eichler integral fixed in Section~\ref{sec:completion-shadow},
giving the harmonic Maass completion of \(H\).  GHV construct the
twined character functions for the \(26\) ordinary \(M_{24}\) classes,
and Gannon proves that the resulting class functions \(H_n\) are
genuine \(M_{24}\)-characters for every \(n\geq 1\).
\end{proof}

\section{References}
\label{sec:references}

\begin{itemize}
\item T. Eguchi, H. Ooguri, and Y. Tachikawa, \emph{Notes on the K3
surface and the Mathieu group \(M_{24}\)}, Experimental Mathematics
20 (2011), 91-96.
\item M. R. Gaberdiel, S. Hohenegger, and R. Volpato, \emph{Mathieu
Moonshine in the elliptic genus of K3}, JHEP 10 (2010), 062.
\item T. Eguchi and K. Hikami, \emph{Note on twisted elliptic genus of
K3 surface}, Physics Letters B 694 (2011), 446-455.
\item T. Gannon, \emph{Much ado about Mathieu}, Advances in Mathematics
301 (2016), 322-358.
\item S. Zwegers, \emph{Mock theta functions}, PhD thesis, Utrecht
University, 2002.
\item A. Dabholkar, S. Murthy, and D. Zagier, \emph{Quantum black
holes, wall crossing, and mock modular forms}, 2012.
\item M. Eichler and D. Zagier, \emph{The Theory of Jacobi Forms},
Birkhauser, 1985.
\item V. Gritsenko and V. Nikulin, \emph{Automorphic forms and
Lorentzian Kac-Moody algebras}, International Journal of Mathematics
9 (1998), 153-199.
\item V. Gritsenko, \emph{Elliptic genus of Calabi-Yau manifolds and
Jacobi and Siegel modular forms}, Algebra i Analiz 11 (1999), 100-125.
\item R. Borcherds, \emph{Automorphic forms with singularities on
Grassmannians}, Inventiones Mathematicae 132 (1998), 491-562.
\end{itemize}

\end{document}
